%; whizzy chapter -dvi
% -initex iniptex -latex platex -format platex -bibtex jbibtex -fmt fmt
% 以上 whizzytex を使用する場合の設定。
 
%     Tokyo Debian Meeting resources
%     Copyright (C) 2012 Junichi Uekawa
%     Copyright (C) 2011 Nobuhiro Iwamatsu

%     This program is free software; you can redistribute it and/or modify
%     it under the terms of the GNU General Public License as published by
%     the Free Software Foundation; either version 2 of the License, or
%     (at your option) any later version.

%     This program is distributed in the hope that it will be useful,
%     but WITHOUT ANY WARRANTY; without even the implied warranty of
%     MERCHANTABILITY or FITNESS FOR A PARTICULAR PURPOSE.  See the
%     GNU General Public License for more details.

%     You should have received a copy of the GNU General Public License
%     along with this program; if not, write to the Free Software
%     Foundation, Inc., 51 Franklin St, Fifth Floor, Boston, MA  02110-1301 USA

%  preview (shell-command (concat "evince " (replace-regexp-in-string "tex$" "pdf"(buffer-file-name)) "&"))

%%ここからヘッダ開始。

\documentclass[mingoth,a4paper]{jsarticle}
\usepackage{monthlyreport}
% 日付を定義する、毎月変わります。
\newcommand{\debmtgyear}{2015}
\newcommand{\debmtgmonth}{1}
\newcommand{\debmtgdate}{17}
% started from zero:
% (let ((year 2013) (month 7)) (+ (* (- year 2005) 12) month -1))
\newcommand{\debmtgnumber}{122}

\begin{document}

\begin{titlepage}
\thispagestyle{empty}
% タイトルページ:編集必要な部分は最初のマクロに飛ばすこと

\vspace*{-2cm}
第\debmtgnumber{}回 東京エリア Debian 勉強会資料\\
\hspace*{-2cm}
\includegraphics{image2012-natsu/dotdeb.pdf}\\
\hfill{}\debmtgyear{}年\debmtgmonth{}月\debmtgdate{}日

% ここはアップデートすること
% 全角文字にしないとフォントのサイズが合わないので注意
\rotatebox{10}{\fontsize{30}{30} {\gt 特集：Emacs関連パッケージの作成}}\\

\vspace*{-2cm}
\hfill{}\includegraphics[height=6cm]{image200502/openlogo-nd.eps}
\end{titlepage}

\newpage

\begin{minipage}[b]{0.2\hsize}
 \definecolor{titleback}{gray}{0.9}
 \colorbox{titleback}{\rotatebox{90}{\fontsize{80}{80} {\gt デビアン勉強会} }}
\end{minipage}
\begin{minipage}[b]{0.8\hsize}
\hrule
\vspace{2mm}
\hrule
\begin{multicols}{2}
\tableofcontents
\end{multicols}
\vspace{2mm}
\hrule
\end{minipage}

\dancersection{事前課題}{野島 貴英}

今回の事前課題は以下です:
\begin{enumerate}
\item 本日、何の作業をやるかを宣言ください。
\item (オプション) どこで今回の勉強会の開催を知りましたか？
\item (オプション) 何について聞きたい／参加者と話をしたいですか？
\end{enumerate}
この課題に対して提出いただいた内容は以下です。
\begin{multicols}{2}
{\small
\begin{prework}{ $BLnEg(B }
  \begin{enumerate}
  \item Q.hack time$B$K2?$r$7$^$9$+!)(B\\
    A. RC$B%G%P%C%0!"(BDebian unstable$B>e$NI>2A!"%P%0%l%]=q$-$J$I!#(BRC$B$O(B\url{https://bugs.debian.org/release-critical/}$B;2>H!#(B
  \item ($B%*%W%7%g%s(B)Q.$B2?$K$D$$$FJ9$-$?$$!?;22C<T$HOC$r$7$?$$$G$9$+!)(B\\
    A. $B<+M3%=%U%H%&%'%"$H<+M3$J4D6-$H<+M3$J5;=Q$K$D$$$F$NG.$$;W$$!*<+M3%=%U%H%&%'%"$H<+M3$J4D6-$H5;=Q$G<R2q$X2?$r%"%T!<%k$7$?J}$,$"$J$?$K$H$C$F(BHappy$B$+!"$=$7$F!"M_K>$r!"$$$?$:$i?4$rK~$?$;$k$N$+!#(B
  \end{enumerate}
\end{prework}

\begin{prework}{ alohaug }
  \begin{enumerate}
  \item Q.hack time$B$K2?$r$7$^$9$+!)(B\\
    A. debian livecd for gnuk key generation
  \item ($B%*%W%7%g%s(B)Q.$B$I$3$G:#2s$NJY6/2q$N3+:E$rCN$j$^$7$?$+!)(B\\
    A. $B$=$NB>(B
  \item ($B%*%W%7%g%s(B)Q.$B2?$K$D$$$FJ9$-$?$$!?;22C<T$HOC$r$7$?$$$G$9$+!)(B\\
    A. gnuk$B$dHkL)80J]4I%*!<%W%s%O!<%I$N:r:#(B
  \end{enumerate}
\end{prework}

\begin{prework}{ keiichirou\_miyamoto }
  \begin{enumerate}
  \item Q.hack time$B$K2?$r$7$^$9$+!)(B\\
    A. Linux$B$NJY6/Ey(B
  \end{enumerate}
\end{prework}

\begin{prework}{ henrich }
  \begin{enumerate}
  \item Q.hack time$B$K2?$r$7$^$9$+!)(B\\
    A. $BB?J,869F=q$$$F$^$9!D!#(B
  \item ($B%*%W%7%g%s(B)Q.$B$I$3$G:#2s$NJY6/2q$N3+:E$rCN$j$^$7$?$+!)(B\\
    A. $BM'C#$dCN$j9g$$$+$iD>@\(B
  \item ($B%*%W%7%g%s(B)Q.$B2?$K$D$$$FJ9$-$?$$!?;22C<T$HOC$r$7$?$$$G$9$+!)(B\\
    A. $B3'$,(BDebian$B$K4X$7$FCN$j$?$$$3$H$rGD0.$7$?$$$G$9(B
  \end{enumerate}
\end{prework}

\begin{prework}{ koedoyoshida }
  \begin{enumerate}
  \item Q.hack time$B$K2?$r$7$^$9$+!)(B\\
    A. $B$^$C$?$j%a%s%F%J%s%9$J$I!#(B
  \item ($B%*%W%7%g%s(B)Q.$B$I$3$G:#2s$NJY6/2q$N3+:E$rCN$j$^$7$?$+!)(B\\
    A. Debian JP$B$N%a!<%j%s%0%j%9%H(B
  \end{enumerate}
\end{prework}

\begin{prework}{ dictoss }
  \begin{enumerate}
  \item Q.hack time$B$K2?$r$7$^$9$+!)(B\\
    A. Debian GNU/kFreeBSD$B$r%M%?$K$7$?%W%l%<%s%F!<%7%g%s=`Hw(B
  \item ($B%*%W%7%g%s(B)Q.$B$I$3$G:#2s$NJY6/2q$N3+:E$rCN$j$^$7$?$+!)(B\\
    A. Debian JP$B$N%a!<%j%s%0%j%9%H(B
  \end{enumerate}
\end{prework}

\begin{prework}{ kenhys }
  \begin{enumerate}
  \item Q.hack time$B$K2?$r$7$^$9$+!)(B\\
    A. RC$B%P%0$X$NBP1~(B
  \item ($B%*%W%7%g%s(B)Q.$B$I$3$G:#2s$NJY6/2q$N3+:E$rCN$j$^$7$?$+!)(B\\
    A. Debian JP$B$N%a!<%j%s%0%j%9%H(B
  \end{enumerate}
\end{prework}


}
\end{multicols}

\dancersection{Debian Trivia Quiz}{野島 貴英}

 Debianの昨今の話題についてのQuizです。

今回の出題範囲は\url{debian-devel-announce@lists.debian.org} や \url{debian-news@lists.debian.org}に投稿された
内容などからです。

\begin{multicols}{2}
%; whizzy-master ../debianmeetingresume201311.tex
% $B0J>e$N@_Dj$r$7$F$$$k$?$a!"$3$N%U%!%$%k$G(B M-x whizzytex $B$9$k$H!"(Bwhizzytex$B$,MxMQ$G$-$^$9!#(B
%

\santaku
{12/29$B$K!"(BDebian$B$N%5!<%S%9$HO"7H$G$-$k(Bchromium$B8~$1$N3HD%5!G=$,%j%j!<%9$5$l$^$7$?!#2?$HO"7H$G$-$k$N$G$7$g$&$+!)(B}
{tracker.debian.org}
{sources.debian.net}
{www.debian.or.jp}
{B}
{zack$B$5$s$N(BDebconf14$B$N%W%l%<%s$r8+$F%$%s%9%Q%$%"$5$l$?$H$N$3$H$G!"(BRaphael Geissert$B$5$s$,:n$C$?(Bchromium$BMQ(B(chrome$B$G$bF0$/!K$N%*%s%i%$%s%=!<%9%3!<%I%(%G%#%?$r<B8=$9$k3HD%5!G=$G$9!#(BDebian$B%Q%C%1!<%8$N%=!<%9%3!<%I$rD>@\%*%s%i%$%s$G1\Mw$G$-$k(Bsources.debian.net$B>e$N%3!<%I$r$3$N3HD%5!G=$r;H$C$F%7!<%`%l%9$KJT=82DG=$K$J$j$^$9!#(Bzack$B$5$s$N<($7$??<9o$JLdBj!JBh(B117$B2sEl5~%(%j%"(BDebian$BJY6/2q;qNA;2>H!K$r2r7h$9$k0Y$N2h4|E*$J%"%$%G%"$N#1$D$G$9!*%^%B$9$P$i$7$$!*<+M3%=%U%H%&%'%"%i%V!"(BDebian$B%i%V$J?M$O!"AaB.!"(Bchromium$B$X(BNautilus$B$+$i$3$N3HD%5!G=$rEj$29~$s$G;n$9$Y$7!*(B}

\santaku
{2014/11/11$B$K(BRapha\"el Hertzog$B$5$s$K$h$jDs0F$5$l$?(BDEP14$B$NFbMF$O0J2<$N$I$l!)(B}
{git$B%Q%C%1!<%8%s%0%j%]%8%H%j$N;H$$J}$K$D$$$F(B}
{control$B%U%!%$%k$K(BDebian-Homepage$B%U%#!<%k%I$rF~$l$k7o(B}
{debian/upstream/metadata$B$rMQ0U$9$k7o(B}
{A}
{$B%Q%C%1!<%8%s%0%j%]%8%H%j$N%V%i%s%A$H%?%0$K%Y%s%@!<>pJs!JNc!'(Bdebian/*,ubuntu/*)$B$d$i%P!<%8%g%s>pJs$NF~$lJ}$J$I$r4^$`!"=t!9$N%k!<%k$NDs0F!#>\$7$/$O!"(Bhttp://dep.debian.net/deps/dep14/$B;2>H!#$J$*!"(BDebian-Homepage$B%U%#!<%k%I$N7o$O(BDEP13$B!#(Bdebian/upstream/metadata$B$N7o$O(BDEP12$B!#(B}

\santaku
{2015/1/8 Marvel$B<R$h$j(BDebian ARM$BMQ%Q%C%1!<%8%S%k%I%^%7%s$N4sIU$,$"$j$^$7$?!#$5$F2?Bf$NDs6!$,$"$C$?$G$7$g$&!)(B}
{7$BBf(B}
{8$BBf(B}
{10$BBf(B}
{B}
{Marvel$B<R$,Ds6!$7$?%5!<%P$O(B MV78460 SoC$B%Y!<%9$N%\!<%I(B8$B$D$rDs6!$H$N$3$H!#$J$*!"(BMV78460$B$O!"(BARM v7$B%3%"#44p$rEk:\$9$k(BMarvel$B<R3+H/$N(BSoC$B!#;29M!'(Bhttp://www.marvell.com/embedded-processors/armada-xp/}

\santaku
{2015/1/10$B$K(Bwheezy$B$,%"%C%W%G!<%H$5$l$^$7$?!#$5$F%P!<%8%g%s$O<!$N$I$l!#(B}
{7.6}
{7.7}
{7.8}
{C}
{wheezy$B$r$*;H$$$NJ}!9$OAaB.%"%C%W%0%l!<%I$7$^$7$g$&!#(B}

\santaku
{2015/1/1$B$K$"$kD9$5L$K~$N(BDebian keyring$B$N>C5n$,40N;$7$?$=$&$G$9!#$"$kD9$5$H$O(B?}
{2048}
{4096}
{8192}
{A}
{$B%"%J%&%s%9$H!">C$5$l$?3:Ev$N%-!<0lMw$O(Bdebian-devel-announce$B$N(B2015/1/1$B$N5-;v$K5-:\$,$"$j$^$9!#(Bhttps://lists.debian.org/debian-devel-announce/2015/01/msg00000.html}

\santaku
{2014/12/6$B$K(Bcdn.debian.net$B$N%l%3!<%I$N<($9@h$r!"$H$"$k(BFQDN$B$N(BCNAME$B$H$7$?$$;]$NLd$$9g$o$;$,(Bdebian-devel ML$B$K$"$j$^$7$?!#$I$3$K8~$1$k!)(B}
{ftp.debian.or.jp}
{ftp.jp.debian.org}
{http.debian.net}
{C}
{cdn.debian.net$B$rDs6!$7$F$$$k(BAraki$B$5$s$+$i$NLd$$9g$o$;!#(Bsqueeze$B0J>e$GEk:\$5$l$F$$$k(Bapt ver 0.7.21$B$G$O(Bhttp$B$N(Bredirect$B$KBP1~$7$F$$$k$N$G!"MxMQ<T$K:G$b6a$$%Q%C%1!<%8%@%&%s%m!<%I@h$r65$($F$/$l$k5!9=$r(Bhttp.debian.net$B$K0lK\2=$7$?$$$H$N$3$H!#FC$KH?BP$9$k?M$O5o$J$$LOMM$G$9!#(B}

\santaku
{2015/1/9$B$K(Bdebian-devel ML$B$K$F!"0lIt%Q%C%1!<%8$N(Bdescription$B$,D9$9$.$k7o$K$D$$$F$N5DO@$,$"$j$^$7$?!#:G$bD9$$(Bdescription$B$N9T$r;}$D%Q%C%1!<%8$O<!$N$I$l!)(B}
{firmware-linux-nonfree}
{texlive-latex-extra}
{irssi-scripts}
{B}
{2015/1/9$B8=:_!"(Bunstable/main$B$G$O!":G$bD9$$9T?t$N=g$G$$$/$H!"(B1$B0L(B texlive-latex-extra($B<B$K(B1935$B9T(B!!)$B!"(B2$B0L(B texlive-fonts-extra(437$B9T(B)$B!"(B3$B0L(B irssi-scripts(350$B9T(B)$B!#(Bunstable/non-free$B$G$O!"(B1$B0L(B firmware-linux-nonfree(251$B9T(B)$B$G$9!#(B}



\end{multicols}

\dancersection{最近のDebian関連のミーティング報告}{野島 貴英}

\subsection{第121回東京エリアDebian勉強会}

\begin{itemize}
\item 場所はスクウェア・エニックスさんのセミナルームをお借りしての開催でした。
\item 参加者は9名でした。
\item セミナ内容は 
\begin{itemize}
\item @kenhysさんにより「DebianとFedoraでパッケージをリリースするまでの話」
\item 野島さんにより、「Debianから見たLinux Mint」
\end{itemize}
でした。
\item 残りの時間でhack timeを行い、成果発表をしました。
\item 宴会の代わりに、「大戸屋ごはん処 新宿イーストサイドスクエア」で夕食会をやりました。
\end{itemize} 

　セミナですが、@kenhysさんにより、groongaについての概要とDebian/Fedoraへ公式パッケージをリリースするまでが語られました。現在のDebian/Fedoraの公式パッケージ化に関して、実際のところが語られました。実際に公式パッケージにするまでの諸々について語られる機会は、他ではあまり見ない件なので、非常に貴重な発表だったと思います。また、野島さんにより、Linux Mintについて、Debianとの比較について語られました。

% % (query-replace-regexp "<.*?>" "")
% % (query-replace-regexp "^[	 ]\+" "")

%-------------------------------------------------------------------------------
\dancersection{Emacs関連パッケージのDebianパッケージ作成}{henrich}
%-------------------------------------------------------------------------------
\index{debian-emacs-package}

\subsection{はじめに}

 Debianでelispパッケージをお作法に従って作る方法について語ってみます。

 emacs関連パッケージのDebianパッケージ化の手続きを簡単にいうと、

\begin{description}
\item [Step 1.] emacsen-commonパッケージを入れる
\item [Step 2.] /usr/share/doc/emacsen-common/debian-emacs-policy.gz を読む
\item [Step 3.] ポリシーに従って作る
\end{description}

以上となります。

…というのも突き放しすぎなので、step by stepで作り方を説明してみます。

\subsection{step by stepでの作り方}

 まず、ソースを展開しておいて、dh-makeパッケージを入れてdh\_makeコマンドでパッケージを作ります。パッケージはdh-make、コマンドはdh\_makeと「-」と「\_」の違いがあるのに注意。この際、*.elなファイルがあるとテンプレートファイルをコピーしてパッケージ名に合わせて修正してくれます（あるいは --with-emacs オプションを使います）。以前は特に.elなファイルがなくてもテンプレートファイルがコピーされていたのですが、毎度消すのが鬱陶しかったのでパッチを送って修正してもらいました(\footnote{\url{http://bugs.debian.org/696793}})

\begin{commandline}
$ sudo apt-get install dh-make
$ dh_make --createorig

Type of package: single binary, indep binary, multiple binary, library, kernel module, kernel patch?
 [s/i/m/l/k/n] s

Maintainer name  : Hideki Yamane
Email-Address    : henrich@XXXXXX
Date             : Sun, 04 Jan 2015 12:26:09 +0900
Package Name     : ag-el
Version          : 0.44
License          : blank
Type of Package  : Single
Hit <enter> to confirm: 
Currently there is no top level Makefile. This may require additional tuning.
Done. Please edit the files in the debian/ subdirectory now. You should also
check that the ag.el Makefiles install into $DESTDIR and not in / .

\end{commandline}
\begin{commandline}
$ ls debian/
README.Debian      changelog  emacsen-install.ex  manpage.sgml.ex  preinst.ex
README.source      compat     emacsen-remove.ex   manpage.xml.ex   prerm.ex
ag.el.cron.d.ex    control    emacsen-startup.ex  menu.ex          rules
ag.el.default.ex   copyright  init.d.ex           postinst.ex      source
ag.el.doc-base.EX  docs       manpage.1.ex        postrm.ex        watch.ex
\end{commandline}

 emacsen-*.ex ファイルがあることがわかります。これをリネームして.exポストフィックスを外してやりましょう。余分なファイルも消しておきます。

\begin{commandline}
$ ls
changelog  control    docs             emacsen-remove   rules   watch
compat     copyright  emacsen-install  emacsen-startup  source
\end{commandline}
%$

 スッキリしました。emacsen* ファイルの中身は特にいじる必要はありません。

 次にdebian/controlを修正します。

\begin{commandline}
$ vi debian/control
------debian/controlファイルここから---------------
Package: ag-el
Architecture: any
Depends: ${shlibs:Depends}, ${misc:Depends}
Description: <insert up to 60 chars description>
 <insert long description, indented with spaces>

Package: ag-el
Architecture: all
Depends: ${shlibs:Depends}, ${misc:Depends}, emacsen-common (>= 2.0.8), emacs,
         silversearcher-ag, 
Description: Emacs frontend to ag
 The Silver Searcher (a.k.a. ag) is very fast grep-like program.
 It is faster and has an attractive features than grep.
 ag.el is simple ag frontend for Emacs, loosely based on ack-and-half.el.
------debian/controlファイルここまで---------------
\end{commandline}
%$

 Architecture: は elisp スクリプトなので all に変えます。ポイントはDepends行のemacsen-common ($>=$ 2.0.8), emacsの2つです。

 emacsen-common ($>=$ 2.0.8)の必要性はEmacs Policyに書いてあります。これと併せてdebian/emacs-compatファイルを置きます（中身は単に「0」とだけする）。このファイルはdh\_installemacsenが処理してくれます。man dh\_installemacsenとすると記載があります。

\begin{commandline}
$ man dh_installemacsen
...中略...
FILES
       debian/package.emacsen-compat
           Installed into usr/lib/emacsen-common/packages/compat/package in the package build directory.
...中略...
\end{commandline}
%$
 
 emacsはメタパッケージになっており、Emacs24などの特定バージョンに依存しています。Depends: emacs24などとすると、Emacs25が出た時に再度Depends行を書き換える必要が発生しますが、このメタパッケージを指定すればパッケージ側での都度変更は不要です。

\begin{commandline}
$ apt-cache show emacs
Package: emacs
Source: emacs-defaults
Version: 46.1
Installed-Size: 25
Maintainer: Rob Browning <rlb@defaultvalue.org>
Architecture: all
Depends: emacs24 | emacs24-lucid | emacs24-nox
Description-ja: GNU Emacs エディタ (メタパッケージ)
 GNU Emacs は、拡張可能で自己説明的なテキストエディタです。
 これは常に推奨の最新 Emacs リリースに依存するメタパッケージです。
Description-md5: 21fb7da111336097a2378959f6d6e6a8
Tag: devel::editor, role::dummy, role::metapackage, suite::emacs, suite::gnu,
 use::editing
Section: editors
Priority: optional
Filename: pool/main/e/emacs-defaults/emacs_46.1_all.deb
Size: 1634
MD5sum: 1f115942065ac452467e02377368ee22
SHA1: dbb1343a3d24f60e5038994e3528dd7486e40943
SHA256: c1fad54e790d69b83f32f2612963baba3ea8091ff3ca72c960c7312096223e3a
\end{commandline}
%$

 その他にもXemacsを含むemacsenというパッケージでの指定もありますが、これは仮想パッケージなので単独で指定はできません。指定するとしたら"Depends: emacs | emacsen" のようにします。

 …とここまで書いて「dh-makeでフォローしてくれればいいんじゃねーの？」と気づきました。では、dh-makeパッケージへのパッチを作りましょう…できましたので、BTSを…しました。将来的にはこの部分は「へーそんなのもあるんだー」的になるはず。あ、reportbugで複数ファイルを添付するには1つずつ-Aで指定が必要なようです。

\begin{commandline}
$ debcheckout dh-make
$ cd dh-make
$ git checkout -b support-modern-emacs-policy
...いろいろ変更...
$ git format-patch master
$ ls
0001-add-emacsen-compat-for-modern-Emacs-lisp-package.patch  0002-add-debian-control-file-for-Emacs-add-on.patch 
debian  dh_make  dh_make.1  lib
$ reportbug -A 0001-add-emacsen-compat-for-modern-Emacs-lisp-package.patch \
-A 0002-add-debian-control-file-for-Emacs-add-on.patch dh-make
\end{commandline}

 これでBug\#774545として登録されました。

 しかし、ここで気になることが。dh-makeでのdebian/controlファイルはどうも種類に合わせてコピーされるようです。これをそのまま適用すると、*.elファイルがある場合に否応なしにEmacs add-onパッケージ用のcontrolファイルになるわけで、複数のバイナリパッケージを生成するような場合には嬉しくありません。

 多分、大本のcontrolファイルがあって、Emacs用のが追加される形のほうが嬉しい…のでは？という気がします。…しかし、その場合の「Section: lisp」指定をどうするのか…

\begin{commandline}
Source: #PACKAGE#
Section: lisp
Priority: optional
Maintainer: #USERNAME# <#EMAIL#>
Build-Depends: #BUILD_DEPS#
Standards-Version: #POLICY#
Homepage: <insert the upstream URL, if relevant>
#Vcs-Git: git://anonscm.debian.org/collab-maint/#PACKAGE#.git
#Vcs-Browser: http://anonscm.debian.org/cgit/collab-maint/#PACKAGE#.git/

Package: #PACKAGE#
Architecture: all
Depends: ${misc:Depends}, emacsen-common (>= 2.0.8), emacs | emacsen,
Description: <insert up to 60 chars description>
 <insert long description, indented with spaces>
\end{commandline}
%$

これ、パッケージのところだけを「追記」する形だといいのですかね？

\begin{commandline}
Package: #PACKAGE#
Section: lisp
Architecture: all
Depends: ${misc:Depends}, emacsen-common (>= 2.0.8), emacs | emacsen,
Description: <insert up to 60 chars description>
 <insert long description, indented with spaces>
\end{commandline}
%$

 そうすればソースパッケージの雛形も1箇所でまとめれば済んで、いい感じになるんじゃないか？とか思い始めました。

\begin{commandline}
Source: #PACKAGE#
Section: UNKNOWN
Priority: optional
Maintainer: #USERNAME# <#EMAIL#>
Build-Depends: #BUILD_DEPS#
Standards-Version: #POLICY#
Homepage: <insert the upstream URL, if relevant>
#Vcs-Git: git://anonscm.debian.org/collab-maint/#PACKAGE#.git
#Vcs-Browser: http://anonscm.debian.org/cgit/collab-maint/#PACKAGE#.git/
\end{commandline}

 でもそれは別ハックですね…。

 あとはdebuildしてlintianのwarningなどをちょこちょこと直せば完成です。

\begin{commandline}
$ dpkg --contents /var/cache/pbuilder/result/ag-el_0.44-1_all.deb 
drwxr-xr-x root/root         0 2015-01-04 13:35 ./
drwxr-xr-x root/root         0 2015-01-04 13:35 ./usr/
drwxr-xr-x root/root         0 2015-01-04 13:35 ./usr/share/
drwxr-xr-x root/root         0 2015-01-04 13:35 ./usr/share/doc/
drwxr-xr-x root/root         0 2015-01-04 13:35 ./usr/share/doc/ag-el/
-rw-r--r-- root/root       175 2015-01-04 13:10 ./usr/share/doc/ag-el/changelog.Debian.gz
-rw-r--r-- root/root      1089 2015-01-04 13:35 ./usr/share/doc/ag-el/copyright
-rw-r--r-- root/root      5279 2014-08-05 05:58 ./usr/share/doc/ag-el/README.md.gz
drwxr-xr-x root/root         0 2015-01-04 13:35 ./usr/lib/
drwxr-xr-x root/root         0 2015-01-04 13:35 ./usr/lib/emacsen-common/
drwxr-xr-x root/root         0 2015-01-04 13:35 ./usr/lib/emacsen-common/packages/
drwxr-xr-x root/root         0 2015-01-04 13:35 ./usr/lib/emacsen-common/packages/remove/
-rwxr-xr-x root/root       465 2015-01-04 13:35 ./usr/lib/emacsen-common/packages/remove/ag-el
drwxr-xr-x root/root         0 2015-01-04 13:35 ./usr/lib/emacsen-common/packages/install/
-rwxr-xr-x root/root      1280 2015-01-04 13:35 ./usr/lib/emacsen-common/packages/install/ag-el
drwxr-xr-x root/root         0 2015-01-04 13:35 ./usr/lib/emacsen-common/packages/compat/
-rw-r--r-- root/root         2 2015-01-04 13:35 ./usr/lib/emacsen-common/packages/compat/ag-el
drwxr-xr-x root/root         0 2015-01-04 13:35 ./etc/
drwxr-xr-x root/root         0 2015-01-04 13:35 ./etc/emacs/
drwxr-xr-x root/root         0 2015-01-04 13:35 ./etc/emacs/site-start.d/
-rw-r--r-- root/root      1225 2015-01-04 13:35 ./etc/emacs/site-start.d/50ag-el.el
\end{commandline}
%$

\subsection{まとめ}

\begin{itemize}
\item elispパッケージは emacsen-common パッケージを入れて /usr/share/doc/emacsen-common/debian-emacs-policy.gz ファイルに従ってパッケージを作ろう
\item お手軽にするにはdh-makeパッケージを使おう
\item Policyを読んでtipsを確認する手間を省くにはdh-makeパッケージをハックしてしまえばいい
\end{itemize}

%-------------------------------------------------------------------------------
\dancersection{会場での無線LANのつなぎ方}{野島 貴英,Roger}
%-------------------------------------------------------------------------------
 \subsection{はじめに}

　今回試験として、会場側でフィルタ無しのグローバル回線を用意しました。
ただ、会場側のセキュリティポリシーにより、wpa-psk AES hidden SSIDという
方式での提供となります。

　以下にDebianマシンでの接続方法を記載します。

 また、自分の環境では違うやり方でつながったという方は、野島まで
教えて下さい。こちらでもノウハウとして溜めていく予定です。

 \subsection{wpasupplicant及び/etc/network/interfacesを利用の場合}

 もっとも良いマニュアルは、/usr/share/doc/wpasupplicant/README.Debian.gz
となります。困った場合はこちらも合わせてご参照下さい。

　以下に/etc/network/interfacesの定義について会場の例を記載します。

\begin{commandline}
$ sudo vi /etc/network/interfaces
-----以下のエントリがなければ追記ここから----------
iface wlan0_debian inet dhcp
     wpa-conf /etc/wpa_supplicant/wpa_supplicant_debian.conf
-----以下のエントリがなければ追記ここまで----------
$ sudo vi /etc/wpa_supplicant/wpa_supplicant_debian.conf
-----以下のエントリを追記ここから----------
network={
     ssid=<<会場のSSID>>
     psk=<<会場のパスワード>>
     scan_ssid=1
}
-----以下のエントリを追記ここまで----------
$ sudo chmod 600 /etc/wpa_supplicant/wpa_supplicant_debian.conf
$ sudo ifup wlan0=wlan0_debian
\end{commandline}
%$

 また、ハマってしまった時のデバッグ方法は、
/usr/share/doc/wpasupplicant/README.Debian.gz中の''4. Trubleshooting''の章が便利です。

 \subsection{その他の無線LAN用パッケージを利用の場合}

　すみません、自分が情報を持たないため、現場で教えて下さい。

\cleartooddpage

\vspace*{15cm}
\hrule
\vspace{2mm}
\includegraphics[width=2cm]{image200502/openlogo-nd.eps}
\noindent \Large \bf Debian 勉強会資料\\
\noindent \normalfont \debmtgyear{}年\debmtgmonth{}月\debmtgdate{}日 \hspace{5mm}  初版第1刷発行\\
\noindent \normalfont 東京エリア Debian 勉強会 （編集・印刷・発行）\\
\hrule

\end{document}
